\section{Problema real, contexto, objetivo y pregunta}

\subsection{Contexto}
Los departamentos de emergencias reciben una demanda que cambia con la hora, mientras que
la capacidad suele asignarse por turnos de forma más rígida. Una distribución uniforme de
equipos puede dejar capacidad ociosa en horas tranquilas y formar colas en horas de mayor
llegada. El proyecto evalúa, mediante simulación, si reasignar la \emph{misma} cantidad de
recursos entre bloques horarios reduce el tiempo de espera.

El escenario de referencia es un DE estadounidense descrito por datos públicos de NHAMCS
2022 \citep{nhamcs2022doc}. No pretende caracterizar un hospital en particular ni producir
recomendaciones clínicas.

\subsection{Pregunta de investigación}
\begin{quote}
\emph{¿Cuánto puede reducirse el tiempo de espera de los pacientes al redistribuir una
cantidad fija de equipos de atención entre seis bloques de cuatro horas, sin aumentar las
horas-recurso totales?}
\end{quote}

\subsection{Objetivos}
\textbf{General.} Diseñar, implementar y evaluar una simulación de eventos discretos que
compare una asignación base con una asignación optimizada de recursos.

\textbf{Específicos.}
\begin{enumerate}
  \item Estimar perfiles horarios de llegada y de triaje con datos públicos NHAMCS 2022.
  \item Construir un modelo parametrizable con llegadas aleatorias, prioridades y capacidad
  variable.
  \item Implementar un algoritmo genético que redistribuya recursos sin aumentar el total.
  \item Verificar automáticamente que toda solución propuesta cumpla las restricciones.
  \item Comparar escenarios con métricas comunes, semillas pareadas y una referencia real.
\end{enumerate}

\subsection{Métricas}
La métrica principal es la \textbf{espera media hasta el inicio de la atención}, en
minutos, con sentido de minimización. Las métricas secundarias son la mediana, el
percentil 90, la proporción de pacientes con espera superior a 60 minutos, la cola media y
máxima (pacientes) y la utilización de los equipos (proporción).

\subsection{Alcance y supuestos}
\begin{itemize}
  \item Horizonte de 24 horas en seis bloques de cuatro horas; la capacidad solo cambia al
  inicio de cada bloque.
  \item Cada paciente usa un solo equipo a la vez; el nivel de triaje no cambia durante la
  espera.
  \item La cola respeta prioridad clínica y, dentro de cada prioridad, el orden de llegada.
  \item El servicio es una abstracción de ocupación del equipo, no tiempo continuo de
  contacto con un médico.
  \item Quedan fuera: diagnóstico y resultados clínicos, traslados, disponibilidad
  individual de personal, camas de hospitalización y cualquier predicción para Guatemala.
\end{itemize}
